% Lösungen Einheit 2 von 4 – Scheitelpunktform (zusammenbau v0.1)
\einheitenkopf{Einheit 2 von 4 · Scheitelpunktform}

\erg{13}{a) Scheitel bei $x = 4$ \quad b) $S(2|5)$ \quad c) $S(-4|2)$ \quad d) $S(-2|3)$ \quad e) Scheitel $S(2|1)$, dann $1$ nach rechts $1$ nach oben, $2$ nach rechts $4$ nach oben, nach links genauso: $(1|2)$, $(3|2)$, $(0|5)$, $(4|5)$ \quad f) $f(x) = (x - 5)^2 + 2$ \quad g) $S(2|{-3})$, nach oben geöffnet: $f(x) = (x - 2)^2 - 3$ \quad h) $g(x) = (x - 1)^2 - 2$, Scheitel $S(1|{-2})$ \quad i) $g(x) = -(x - 5)^2 - 1$, Scheitel $S(5|{-1})$ \quad j) $g(x) = (x + 4)^2 + 2$, Scheitel $S(-4|2)$ \quad k) z. B. $f(x) = (x - 4)^2$ mit $S(4|0)$; jede Parabel $(x - d)^2$ mit positivem $d$ \quad l) $g$ ist $f$ um $6$ nach unten verschoben – kein gemeinsamer Punkt: Scheitel $(1|2)$ und $(1|{-4})$ übereinander, gleiche Öffnung. \quad m) $S(4|1)$, nach oben geöffnet, schmaler; Punkt $(5|3)$ \quad n) nach oben: $(x - 4)^2$; gespiegelt: $g(x) = -(x - 4)^2$, Scheitel $S(4|0)$ \quad o) $g(x) = (x - 4)^2 - 2$: Normalparabel mit Scheitel $S(4|{-2})$, Öffnung bleibt \quad p) z. B. $f(x) = -x^2 - 2$ mit Scheitel $(0|{-2})$ (jeder Scheitel unter der $x$-Achse); $g(x) = x^2 + 2$}
\erg{14}{a) In der Klammer steht $+6$, der Scheitel liegt also bei $x = -6$: $S(-6|{-2})$ \quad b) $(x - 4)^2$ ist ein Quadrat und nie negativ; null wird es nur bei $x = 4$ – dort ist $f(x)$ am kleinsten. \quad c) falsch; wahr; falsch – $f(x) = (x - 1)^2 - 4$, $S(1|{-4})$, nach oben geöffnet \quad d) Scheitel $S(6|3{,}8)$: höchstens $3{,}8\,$m hoch, $6\,$m vom Abstoß entfernt}
